\chapter{A new complete R1 theorem with no source-ratio split}
\label{component:pbaab204a8ab1}
\noindent\textbf{Component }\path{pbaab204a8ab1}\par
{\small\noindent Source: \path{NEW/weak_c2/r1_all_sources_all_tests.tex}\par}
{\small\noindent SHA256: \path{baab204a8ab1db9336f65126c06bd3c57110f45d4ea11df2ab591f61d25a4ac1}\par}
{\small\noindent Inventory use: R1N1-R1N24 and all-testing restoration | Only R1N12-24 universal full-current/finite-cap algebra is reused; no original R1 source membership or local theorem is imported\par}
\medskip
% New proof, 2026-10-22.  Frozen reconstruction files are not modified.
\section{A new complete R1 theorem with no source-ratio split}
\label{pbaab204a8ab1:sec:new-r1-all-tests}

\paragraph{Statement and exact source.}
Use the original common-II source, with $d=1$, $\lambda=\beta=0$,
the six native coefficients of the recovered
\path{weak_c2_source_formulas.tex}, every original strict gap,
and $n\geq0$.  In the original R1 sector,
\[
 t_s,\quad B_s,\quad a_sB_s,\quad t_sB_s>0
 \quad\hbox{are infinitesimal}.                         \tag{R1N1}
\]
There is no bound on $a_s$, on any original positive opening, or on
the source ratio previously called $q_{\rm old}$.  The conclusion is:
the entire original parent is a finite separated tensor sum of
squares at every independent real projective measurement pair and
every real or complex state.  This proof treats R1 only.

Source and testing quantities belong to separate real closed fields
$K_c,K_t\supset\mathbb R$.  All SOS assertions are in
$K_c\otimes_{\mathbb R}K_t$.  We use the actual recovered strict
bounded-Gram lemma: a bounded finite Gram with positive definite
real residue gives tensor SOS, with no boundedness requirement on
its row entries.  No order or mixed inverse is imposed on the
tensor ring.

Put $c_s=1+a_sB_sr$, $R_s=1+ha_sB_s$, and define
\[
\begin{aligned}
 j&=\frac{a_sc_s}{R_s},\quad B=\frac{B_sR_s}{c_s},\quad
 T=\frac{t_sc_s}{R_s},\quad p=\frac{a_s(1-r)}{R_s},\\
 H&=\frac{1-h}{t_sc_s},\qquad x=jB,\quad y=TB,\quad q_z=jT.
\end{aligned}                                          \tag{R1N2}
\]
The native strict gaps give $0<p<\min(1,j)$ and $0<H<1$.
They also give $1-h<t_sc_s$, so $h\sim1$.
Consequently $B,T,x,y$ are infinitesimal, while
\emph{$q_z$ is an arbitrary positive source quantity}.
In fact $x=a_sB_s$ and $y=t_sB_s$ exactly.
The inverse source substitution is
\[
 r=\frac{j-p}{j+px},\quad h=\frac{1-HT}{1+HTx},\quad
 a_s=\frac{j+px}{1+HTx},\quad
 B_s=\frac{B(1+HTx)}{1+pB},\quad
 t_s=\frac{T(1+pB)}{1+HTx}.                             \tag{R1N3}
\]
All arguments below therefore concern the actual original source.

\subsection{Source bounds without a bound on \(jT\)}

The exact dictionary RSR4--RSR6 in the freshly reconstructed
actual-source and scalar proof remains valid under
(R1N2)--(R1N3), because it is an exact rational substitution.
For completeness the two polynomials controlling its denominators
and current masses are
\begin{align*}
 \mathcal P={}&j+1+B\{(j+1)^2-p\}\\
 &+TB\{Hj+j+1+B[Hj(2j+2-p)+p(2j+1-p)]\}\\
 &+HjT^2B^2(1+Bp+BHj),\\
 \mathcal L={}&K+T\{(1-H)j+1-p-H\}\\
 &+BT[pK+Hj(1-p)]+BHT^2\{(1-H)j-p\},
 \qquad K=j+1-p .
\end{align*}
The estimates
\[
 \mathcal P\sim j+1,\qquad \mathcal L\sim K,\qquad
 \frac12<k:=\frac K{j+1}<1                               \tag{R1N4}
\]
use only $T,B,Bj=x$ infinitesimal and $p,H<1$.
Indeed dividing each correction in $\mathcal P$ by $j+1$ bounds it
by a finite sum of $B+x$, $TB(1+B+x)$ and
$T^2B^2(1+B+x)$; dividing $\mathcal L-K$ by $K$ bounds it
by a finite sum of $T$, $TB$ and $BT^2$.
In particular no product $jT$ is bounded in this argument.

The exact nonnegative numerator of $n$ has the regrouping
\[
 \mathcal N=p\{1+B(3j+1-p)+BT(j+1)\}
       -Bj(j+1)\{1+(1-H)T\}+\mathcal R ,
 \quad |\mathcal R|\le C_0BT\,Bj(j+1).                   \tag{R1N5}
\]
This is the finite polynomial regrouping RSR8; its coefficient bound
uses $p<\min(1,j)$ only.  Since
$n=\mathcal N/[2(j-p)\mathcal P]\ge0$, it follows that
\[
 \eta_*:=\frac{Bj(j+1)}p=O(1),\qquad
 \eta_*\in\mathrm I
 \Longrightarrow n\sim\frac{p}{2(j-p)(j+1)}>0.           \tag{R1N6}
\]
Thus literal $n=0$ is retained.  The same dictionary gives, with
$\mu=m(1-p)$,
\[
\begin{aligned}
 \mu&=O(p),\quad \mu+\tfrac12\asymp1,\quad
 \alpha\asymp B^2(j+1),\quad \frac{n(j-p)}j=O(1),\\
 \frac AT&\sim\frac k2,\quad
 TCH\sim\frac k2,\quad TD(1-H)\sim\frac k2.
\end{aligned}                                          \tag{R1N7}
\]
All statements are uniform at the independent openings $H$, $1-H$,
$j-p$, and $1-p$.

Put
\[
 g_*=\frac{CD}{C+D},\qquad h_*=\frac D{C+D}.
\]
Writing $U_c=TCH$ and $U_D=TD(1-H)$ gives the exact source identity
\[
 Tg_*=\frac{U_cU_D}{(1-H)U_c+HU_D}\sim\frac k2.          \tag{R1N8}
\]
The denominator is a bounded positive unit even when either opening
is infinitesimal.  In particular
$1/C,1/D,1/g_*$ are infinitesimal and
$Ag_*\sim k^2/4$.  Also $(C+D)^{-1/2}=O(\sqrt T)$.

\subsection{The complete small-angle theorem extends to all R1 sources}

The reconstructed proof \path{fn_reconstructed.tex}, RFN4--RFN24,
is a full-parent proof, not just a restriction to a scalar graph.
We now justify its application under the larger source domain
(R1N1), rather than importing its former domain label.

First, all of its exact identities hold for the rational dictionary
(R1N3), without any size assumption on $q_z$.  The parent, its two
normals, their quadratic block, the scalar circuit, the radial Gram
and the three source cancellation identities were all proved as
identities of this dictionary.  Second, every uniform source bound
used in that proof is exactly (R1N4)--(R1N7) or the following scalar
bound:
\[
 f_c\asymp Tx^2,\quad q_c\asymp Tx,\quad
 r_c\asymp TjB^2(1-p),\quad
 g_c/f_c=O(Hy),\quad |k_c|/(q_cy)=O(p+B).                \tag{R1N9}
\]
To check the latter bounds without $jT\in\mathrm I$, use the explicit
finite polynomials RSR15 and RSR18: each term of
$\mathcal F_1/K,\mathcal F_2/K$ is bounded using $p/j<1$
and $Bj=x$, and $\mathcal K/K=O(p+B)$ by the same inequalities.
Thus $\mathcal F\sim K$ and all statements (R1N9) follow from
RSR16--RSR20.  No $jT$ factor occurs in these estimates.

The signed scalar allocation retains the effective mass
$M_e=nh_0j^2/(n+h_0)$, $h_0=\mu+1/2$, after the complete
correlated $n$-square completion.  Its three source cases use
$1-p$ a unit, or $p\sim1$ with $\eta_*$ a unit, or
$p\sim1$ with $\eta_*$ infinitesimal.  The proofs of these cases use
only (R1N4)--(R1N9).  In the last case $M_e$ is a unit; in the second,
$j^2B$ is a unit.  Literal $n=0$ occurs only in the first two cases.

For clarity, the following are all potentially nonuniform costs in
the complete normal lift, after cancellation of the source factors.
The symbols $u,v$ denote independent infinitesimal tests.
\[
\begin{array}{c|c}
 \text{purchased row or source case}&\text{bounded diagonal costs}\\ \hline
 v(u^2-x),\ 1-p\text{ unit}
       &O(B),\ O(Bu^2),\ O(u^2/(1-p))\\
 v(u^2-x),\ p\sim1,\ \eta_*\text{ unit}
       &O(B+u^2/j+u^4/(j^2B)),\ O(Bu^2)\\
 v(u^2-x),\ p\sim1,\ \eta_*\in\mathrm I
       &O(u^2/j),\ O(u^4/j),\
            [8u^4+2(u^2-x)^2]/M_e\\
 pv^2(v-uz/j),\ \eta_*\in\mathrm I
       &(p^2/M_e)v^4\\
 pv^2(v-uz/j),\ \eta_*\text{ unit}
       &(p^2y/f_c)(yv^2+v^4)
\end{array}                                                   \tag{R1N10}
\]
Every cost is infinitesimal for its stated case.  The source quotients
$p^2/M_e$ or $p^2y/f_c$ used in the last two lines are bounded by
(R1N6) and $p<\min(1,j)$.

The full opposite normal row has squared multiplier
\[
 a_R=Tq_zv^2(\alpha-BTCHu^2)^2,\qquad
 J=B(j+1)+u^2 .
\]
The occurrence of the possibly unbounded $q_z$ cancels exactly in
its three cost ratios:
\[
 a_R/f_c=O((T/j)v^2J^2),\qquad
 a_R/(q_cy)=O(v^2J^2),\qquad
 a_R/(M_eyx)=O\!\left(T(1+1/n)v^2J^2/j^2\right).        \tag{R1N11}
\]
The last ratio is used only for $p\sim1,\eta_*\in\mathrm I$,
where $(1+1/n)/j^2$ is bounded.  In the other two cases the
direct opposite-current or the $(1+U^2)\Psi$ budget applies exactly
as in RFN24.  The remaining normal coefficients are bounded
multiples of $uv\sqrt T$, $T+v^2$ or the rows in (R1N10);
their estimates require only $B,T,x,y$ infinitesimal.
This checks all complete first-normal and hard-normal coefficients.
The variable $z$ remains inside retained rows, with no bound on it.

Consequently RFN4--RFN24 proves the following enlarged local statement:
for every original R1 source (R1N1), the entire parent is tensor SOS
when $u,v>0$ are infinitesimal, $w=(W+uY)/(1+u^2)=1$, and
$z=(Z-uX)/(1+u^2)>0$ is arbitrary.  The complete source-only
$n$-completion includes $n=0$.  The freshly proved SJ1--SJ10
state/sign reduction therefore gives every state for these small
angles, as well as every literal zero-angle and projective endpoint.

\subsection{Outer testing angles: two exact current identities}

We include the outer mechanism and its precise positivity condition.
Let
\[
 R=uW-Y,\quad C_t=v(uX-Z),\quad B_t=X+uZ,\quad
 N=R+h_*C_t,\quad D_u=1+u^2,\quad D_v=1+v^2,
\]
and let $K$ be the signed quarter-core
\[
 K=\tfrac12(R+C_t)^2+\tfrac12D_u(vY-Z)^2
             +\tfrac12D_v B_t^2-\tfrac12D_uD_vZ^2.       \tag{R1N12}
\]
All unselected original nonnegative events are kept as whole squares.

If $1/v$ is bounded, use
\[
 K_0=2\sqrt{g_*}C_t,\quad
 B_0=\sqrt{2D_v/D_u}\,B_t,\quad
 c_u=D_u^{-1/2},\quad t_u=uD_u^{-1/2},\quad
 \eta_v=\sqrt{D_v}/v .
\]
The exact full selected-parent identity is
\begin{align*}
 4CD_vR^2+4D(R+C_t)^2+4K
 ={}&4(C+D)N^2+4Cv^2R^2+2(R+C_t)^2\\
 &+2D_u(vY-Z)^2+(K_0,B_0)Q(K_0,B_0)^T,\\
 Q={}&
 \begin{pmatrix}
 1-c_u^2\eta_v^2/(2g_*)&t_u\eta_v/\sqrt{2g_*}\\
 t_u\eta_v/\sqrt{2g_*}&1
 \end{pmatrix}.
\tag{R1N13}
\end{align*}
Since $1/g_*$ is infinitesimal and all the testing factors are bounded,
$\overline Q=I$.  This treats every other angle and every state,
including Bob infinity by homogeneous clearing.

For $v$ infinitesimal put $\epsilon=1/g_*$, $q=\sqrt\epsilon$,
$\gamma=Ag_*$, and
\[
 K_1=2\sqrt{g_*}C_t+\sqrt\epsilon Y,\quad
 B_1=B_t-uvY,\quad V_0=vY,\quad Y_1=\sqrt{\epsilon D_u}\,Y,
\]
\[
 L_0=\tfrac12qc_u(K_1-c_uY_1)-t_uvB_1+c_uvV_0.
\]
Then $L_0=\sqrt{D_u}\,v(vY-Z)$ exactly, and
\begin{align*}
 &4CD_vR^2+4D(R+C_t)^2+4AD_uD_vY^2+4K\\
 &\quad=4(C+D)N^2+4Cv^2R^2+2(R+C_t)^2+\mathcal G,\\
 \mathcal G={}&K_1^2+2D_v(B_1^2+V_0^2)
 +(4\gamma D_v-c_u^2)Y_1^2\\
 &+2qv(K_1-c_uY_1)V_0-2L_0^2 .
\tag{R1N14}
\end{align*}
The complete four-root Gram has residue
\[
 \operatorname{diag}(1,2,2,\overline{k}^{\,2}
                                      -\overline{c_u}^{\,2}).    \tag{R1N15}
\]
Thus it is strict whenever the last entry is positive.  In particular
it covers every unbounded $u$, and covers every noninfinitesimal
$u$ whenever $\overline k=1$.  All coefficients in the identity
are bounded functions of the displayed normalized homogeneous Alice
coordinates, so it includes the literal Alice-infinite endpoint.

\subsection{The remaining bounded nonzero Alice chart}

The only outer chart left by (R1N13)--(R1N15) has $v$ infinitesimal,
$u$ a bounded positive unit, and $\overline k<1$.  Since
$1-k=p/(j+1)$, this implies that $p,j$ are positive units and $j$
is bounded.  Therefore $\eta_*\in\mathrm I$ and (R1N6) gives the
exact relative critical limits
\[
 m\sim\frac{jp}{2(1-p)(j+1)},\qquad
 n\sim\frac{p}{2(j-p)(j+1)}.                              \tag{R1N16}
\]
These statements allow either or both coefficients to be unbounded.

Choose bounded positive real caps $m_c\le m$, $n_c\le n$ as specified
below and keep their nonnegative excesses.  Retain all of $\alpha$
as its original square.  Define six rows
\[
 \widehat N=\sqrt{C+D}(R+h_*C_t),\quad
 P=C_t/\sqrt T,\quad Y_0=\sqrt T\,Y,\quad
 X,\quad Z,\quad V=vY.                                  \tag{R1N17}
\]
They need not be independent or bounded.  Put
$\delta=(C+D)^{-1/2}=O(\sqrt T)$, so
$R=\delta\widehat N-h_*\sqrt T\,P$.
The selected \emph{entire} parent is exactly
\begin{align*}
 \mathcal P_{\rm sel}={}&
 \widehat N^2+Cv^2R^2+\tfrac12(R+C_t)^2
       +(Tg_*)P^2+(A/T)D_uD_vY_0^2\\
 &+m_c\{(v/u)R+(D_u/u)V-X-uZ\}^2+n_cD_uD_vX^2\\
 &+\tfrac12D_u(V-Z)^2+\tfrac12D_v(X+uZ)^2
                                      -\tfrac12D_uD_vZ^2 .
\tag{R1N18}
\end{align*}
This is obtained by the exact completion
$CD_vR^2+D(R+C_t)^2=(C+D)N^2+g_*C_t^2+Cv^2R^2$.
It keeps every state-normal term.

The row relation
\[
 PY_0=V(uX-Z)                                            \tag{R1N19}
\]
allows an arbitrary real zero shift
$2b(PY_0-V(uX-Z))$.  Keep the two explicitly nonnegative squares
$Cv^2R^2$ and $(R+C_t)^2/2$ whole.  The remaining complete Gram
on (R1N17) is bounded: $u,1/u$ are bounded, $Tg_*,A/T$ are
units, and the $R$ contribution in the $m_c$ square has coefficients
$\delta v/u$ and $h_*\sqrt T v/u$, both infinitesimal.
Its residue is the direct sum of a unit pivot and the two blocks
\[
 \begin{pmatrix}k_0/2&b\\b&k_0(1+u_0^2)/2\end{pmatrix},\qquad
 Q_3(m_c,n_c;b),                                        \tag{R1N20}
\]
where subscripts denote real residues and $Q_3$ is the quadratic
form in $(X,Z,V)$
\begin{align*}
 Q_3={}&m_c\{(D_u/u)V-X-uZ\}^2+n_cD_uX^2\\
       &+\tfrac12D_uV^2-D_uVZ+\tfrac12(X+uZ)^2
                                            -2bV(uX-Z).
\tag{R1N21}
\end{align*}
In (R1N21), $u=u_0>0$ is real.  The choices of $b,m_c,n_c$
that make both blocks strict are now explicit.

\paragraph{Interior openings.}
Suppose $p_0<1$ and $j_0>p_0$.  Use initially
$m_0=j_0p_0/[2(1-p_0)(j_0+1)]$ and
$n_0=p_0/[2(j_0-p_0)(j_0+1)]$, and take $b=k_0/2$.
The first leading principal minor of $Q_3$ is
$m_0+n_0D_u+1/2>0$; the other two are
\[
 \frac{p_0u_0^2D_uK_0}
      {4(j_0-p_0)(j_0+1)^2(1-p_0)},\qquad
 \frac{p_0u_0^2D_uK_0^2}{8(j_0-p_0)(j_0+1)^3},          \tag{R1N22}
\]
where $K_0=j_0+1-p_0$.  They are positive.
The first block in (R1N20) has determinant
$k_0^2u_0^2/4>0$.  Positive definiteness is open, so fixed
real caps slightly below $m_0,n_0$ preserve strictness and are
smaller than the actual source coefficients.

\paragraph{Only the \(n\) opening is infinite.}
If $j_0=p_0<1$, then $n$ is unbounded and
$m_0=j_0^2/[2(1-j_0)(j_0+1)]$ is finite positive.
At $b=k_0/2=1/[2(j_0+1)]$, restrict (R1N21) to the kernel
$X=0$ of the $n$ packet.  Its first diagonal on $(Z,V)$ is
$u_0^2(m_0+1/2)>0$, and its determinant is
$u_0^2k_0^2/4>0$.  Choose $m_c<m_0$ sufficiently close and
then a sufficiently large fixed real $n_c$.  The whole block is
strict, with both caps below the original weights.

\paragraph{Only the \(m\) opening is infinite.}
If $p_0=1<j_0$, then $m$ is unbounded and
$n_0=1/[2(j_0^2-1)]$ is finite positive.
Take a real $b$ just above $k_0/2=j_0/[2(j_0+1)]$, with
\[
 b<\min\{k_0\sqrt{D_u}/2,\ j_0/[2(j_0-1)]\}.
\]
This interval is nonempty.  The first block in (R1N20) is strict.
On the kernel $(D_u/u)V-X-uZ=0$ of the $m$ packet, use $(X,V)$
coordinates.  The first diagonal of (R1N21) is
$D_u/[2(j_0^2-1)]>0$, and its determinant is
\[
 -\frac{D_u^2[\,2b(j_0-1)-j_0\,][\,2b(j_0+1)-j_0\,]}
        {4u_0^2(j_0^2-1)}>0.                            \tag{R1N23}
\]
Choose $n_c<n_0$ sufficiently close and then a sufficiently large
real $m_c$.  This gives the required full strict block.

\paragraph{Both openings are infinite.}
If $p_0=j_0=1$, both $m,n$ are unbounded and $k_0=1/2$.
Choose
$1/4<b<\sqrt{D_u}/4$.
The common kernel of the two event packets is
$X=0$, $Z=D_uV/u_0^2$.  On it the rest of (R1N21) is exactly
\[
 \frac{D_u(4b-1)}{2u_0^2}V^2>0.                        \tag{R1N24}
\]
A sufficiently large real cap on both packets therefore gives
a strict block.

The elementary finite-cap fact used in these cases is as follows:
if a real positive semidefinite event Gram is positive definite on
the complement of its kernel, and the remaining real quadratic
form is positive definite on that kernel, a sufficiently large
real multiple of the event Gram makes their sum positive definite.
It follows by a block Schur complement, and works for the sum of
two event Grams as well.  Thus the choices above prove strictness
of the \emph{complete} bounded Gram (R1N18), without restricting
any normal or state coordinate.

\subsection{Exhaustive testing and state join}

The sign gauges SJ3 preserve every original source packet and the
whole core, so affine $u,v$ may be taken nonnegative.
If $v$ is not infinitesimal, (R1N13) applies.
If $v$ is infinitesimal and $u$ is not infinitesimal, either
(R1N15) applies or the complete six-root argument
(R1N16)--(R1N24) applies.  These cases exhaust the outer angles.
If $u,v>0$ are both infinitesimal, the enlarged local FN theorem
proved above applies after SJ's full-state reduction to $w=1,z>0$;
SJ gives the other state signs and the axes directly.
Literal homogeneous measurement endpoints are given by SJ10 or
the homogeneous current identities, without a limiting argument.
The zero state is immediate.  For a complex state, the real
parent matrix splits its value into its real and imaginary
quadratic evaluations, to which the same proof applies.

Restore every capped excess, the original $n,m,\alpha$ packets
not used in a particular chart, and all cleared positive pure
testing homogeneous factors.  All normalizations, radicals and
inverses are pure source or pure testing.  This proves the
complete original R1 source row at all independent tests and
states, without a source-ratio split or a comparison of source
and testing rates.

\paragraph{Verification scope.}
The companion checker verifies the exact source-coordinate map,
the harmonic source identity, both full outer-current identities,
the new full six-root parent representation, the zero relation,
all positive-minor and kernel formulas, and their original
source coefficients.  The source bounds and exhaustive chart
argument above are analytical.  This theorem uses actual complete
source and FN proof texts, not a historical status receipt.
It makes no assertion about the still-separate R2--R4 rows.
